% ================= 4. DESIGN =================
\section{Design (Preliminary)}

\subsection{System Architecture}

\subsection{Technical Requirements}
\subsubsection{Coordinate Frames}

\subsection{Mathematical Modelling}

\subsubsection{2D model}

A two-dimensional (2D) prototype will be built at the scale of a single face of the cube, in order to develop and validate the control algorithms before the full cube is assembled. The prototype will consist of a 3D-printed plate matching the dimensions of one cube face, with a momentum wheel and motor mounted at its center and a braking mechanism at one corner. The plate will rest on a bearing at that corner, giving it a single rotational degree of freedom within a vertical plane.

The nonlinear equations of motion governing the system will be
\begin{align}
\ddot{\theta}_b &= \frac{(m_b l_b + m_w l)\, g \sin\theta_b - T_m - C_b \dot{\theta}_b + C_w \dot{\theta}_w}{I_b + m_w l^2}, \\
\ddot{\theta}_w &= \frac{(I_b + I_w + m_w l^2)(T_m - C_w \dot{\theta}_w)}{I_w (I_b + m_w l^2)} \notag \\
&\quad - \frac{(m_b l_b + m_w l)\, g \sin\theta_b - C_b \dot{\theta}_b}{I_b + m_w l^2},
\end{align}
with the motor current-torque relation $T_m = K_m u$, where $\theta_b$ is the tilt angle of the pendulum body, $\theta_w$ the rotational displacement of the momentum wheel relative to the body, $m_b$ and $m_w$ the masses of the body and wheel, $I_b$ and $I_w$ their moments of inertia about the pivot and motor axis respectively, $l$ the distance from the motor axis to the pivot, $l_b$ the distance from the body's center of mass to the pivot, $g$ the gravitational acceleration, $T_m$ the motor torque, $C_b$ and $C_w$ the dynamic friction coefficients of the body and wheel, $K_m$ the motor torque constant of the selected motor, and $u$ the motor current input commanded through the motor driver.

Since the pendulum body and wheel will be 3D-printed, their geometry will be fully defined in CAD prior to fabrication. This will allow $m_b$, $m_w$, $I_b$, $I_w$, and $l_b$ to be computed directly from the CAD model's mass-properties simulation, given the chosen print material, infill density, and the mass and placement of mounted hardware such as the motor driver, encoder, microcontroller, IMU, and wiring. These simulated values will be cross-checked experimentally once the prototype is fabricated: $l_b$ will be verified by freely hanging the body from different corners, and $I_b$ will be verified by fixing the wheel, hanging the assembly upside down, letting it swing, and fitting the identified model $(I_b+I_w+m_wl^2)\ddot\theta_b = -C_b\dot\theta_b + (m_bl_b+m_wl)g\sin\theta_b$ to the logged tilt data. The damping coefficients $C_b$ and $C_w$ cannot be obtained from simulation, since they depend on bearing friction, wiring drag, and motor cogging not captured in CAD, and will therefore be identified purely experimentally: $C_w$ (together with $I_w$, as a check against the simulated value) by driving the momentum wheel with current steps while the body is clamped and fitting $I_w\ddot\theta_w = K_mu - C_w\dot\theta_w$ to the logged wheel speed; $C_b$ from the swing-test data described above.

\begin{table}[h]
\centering
\caption{Prototype parameters and determination method}
\begin{tabular}{lll}
\hline
Symbol & Description & Method \\
\hline
$l$ & Motor axis to pivot distance & design choice \\
$l_b$ & Body CoM to pivot distance & design choice \\
$m_b$ & Pendulum body mass & calculated based on design \\
$m_w$ & Momentum wheel mass & calculated based on design \\
$I_b$ & Body moment of inertia & CAD simulation \\
$I_w$ & Wheel moment of inertia & CAD simulation \\
$C_b$ & Body damping coefficient & estimated \\
$C_w$ & Wheel damping coefficient & estimated \\
$K_m$ & Motor torque constant & Motor datasheet \\
\hline
\end{tabular}
\end{table}

With $I_b$, $I_w$, $m_b$, $l_b$, $m_w$, and $l$ established from CAD simulation, the wheel angular velocity $\omega_w$ required for the jump-up maneuver will be estimated before fabrication. Assuming a perfectly inelastic collision between wheel and body, conservation of angular momentum during impact together with conservation of energy from the moment of impact to the edge-standing position gives
\begin{equation}
\omega_w^2 = (2-\sqrt2)\,\frac{I_w+I_b+m_wl^2}{I_w^2}\,(m_bl_b+m_wl)\,g.
\end{equation}
This simulation-derived value of $\omega_w$ will be used to verify, ahead of fabrication, that the selected motor can reach the required speed within the available bus voltage, and to size the wheel's inertia — via its radius and any added rim mass — so that the required speed remains within the motor's practical operating range without excessively reducing the recovery angle available for balancing. Once the prototype is built, the current-step and swing-test identification described above will be used to refine $\omega_w$ experimentally, together with an iterative trial-and-error adjustment to account for any deviation from the idealized inelastic-collision assumption (e.g., a non-zero coefficient of restitution or parameter uncertainty).

The braking mechanism will consist of a servo-driven barrier that is swung into a bolt head attached to the momentum wheel to stop it abruptly. The servo's response time will be a design constraint on the validity of the "instantaneous stop" assumption underlying the jump-up equation above; a slower stop than assumed will yield a lower effective post-impact body velocity than predicted, which will be compensated for through the same iterative tuning process rather than relying on the analytical estimate alone.

The electronics and software will be built around a microcontroller communicating over a CAN bus with a motor driver that closes the current loop for the selected brushless motor. Wheel speed will be measured with a magnetic encoder on the motor shaft, and body tilt will be estimated from an onboard IMU. Braking will be actuated by a digital servo driven directly from the microcontroller, and a wireless module will provide a telemetry and debugging link. Power will be supplied by a LiPo battery, stepped down through a DC-DC converter for the logic-level electronics, with the motor driver powered directly from the battery bus.
\subsubsection{3D model}



\subsubsection{Equations of Motion --- Reaction-Wheel Inverted Pendulum}

\subsubsection{Equilibria and Linearization}

\subsubsection{Control-Oriented Model Summary}

\subsection{Control Architecture}

\subsubsection{LQR Baseline Control}

\subsubsection{Gain Scheduling \& Mode Transitions}

\subsubsection{Jump-Up Strategy}

\subsection{Mechanical Design Intent}

\subsubsection{Structural Concept}

\subsection{Electrical Design Intent}

\subsubsection{Components}

\subsubsection{Electrical Diagram}
